% Round 20. HM20. family20. Depends on final DP20 activation and independent review.
\section{All surface Hilbert schemes through actual Chow transfer}
\label{r20hm:section}

\paragraph{Historical and logical status.}
The R9 source \texttt{root\_chow\_envelope\_and\_hilbert\_schemes.tex}
already states this application conditionally on its product
candidate.  This section verifies the application using the
present DP20 product theorem and the frozen CV16 obstruction-kernel
covariance.  It does not change the old source's candidate status
or claim bibliographic novelty.  Its activation depends on the
final independent acceptance of DP20 and of this transfer.

\subsection{The exact previously proved transfer mechanism}

For a smooth projective $V/K$, let
$\mathcal K_K^r(V)=\ker\mathfrak o_{V,K}^r$ be the frozen
intrinsic obstruction kernel.  CV16 proves that every rational
algebraic correspondence between smooth projective varieties
sends this kernel to the kernel in the degree prescribed by
its codimension.  Its proof uses supported operations,
projective-space integration and the smooth blowup formula;
it is not an assertion of arbitrary proper pushforward of
actual forms.

Suppose smooth projective $X,V_1,\ldots,V_t/K$ have actual
rational Chow correspondences $u_i:X\dashrightarrow V_i$,
$v_i:V_i\dashrightarrow X$ of opposite degree shifts with
\[
 \sum_i v_i\circ u_i=[\Delta_X]
 \quad\text{in }\operatorname{CH}^{\dim X}(X\times X)_{\mathbf Q}.
 \tag{HM20.1}\label{r20hm:diagonal}
\]
If every geometrically homologically trivial class on every
$V_i$ has zero obstruction, the same holds on $X$.
Indeed for $z\in\operatorname{CH}^r(X)_{\mathbf Q,\mathrm{geom\ hom}}$
each $u_{i*}z$ is geometrically homologically trivial, since
cycle classes commute with the actual correspondences.
It belongs to $\mathcal K(V_i)$; CV16 sends $v_{i*}u_{i*}z$
into $\mathcal K(X)$.  Equation~\eqref{r20hm:diagonal}
and additivity give $z\in\mathcal K(X)$.  A term with
impossible Chow degree is zero; a term in geometric
homological codimension zero is also zero, component by
component.  This is the exact diagonal-factorization
corollary CV16.4, restated to identify the required input.

VT17 now constructs the Green form for the original specified
cycle, using its own curved pair and the zero obstruction.
Thus no support inherited from a correspondence image replaces
the original support.  No cohomological or numerical diagonal
identity can substitute for \eqref{r20hm:diagonal}.

\subsection{The native de Cataldo--Migliorini certificate}

Let $S/K$ be a smooth projective geometrically integral surface
over a characteristic-zero field, and put $H=S^{[m]}$, $m\geq1$.
The Hilbert scheme is smooth projective geometrically integral
of dimension $2m$.  For a partition
\[
 \lambda=1^{a_1}\cdots m^{a_m}\vdash m,
 \qquad l=l(\lambda)=\sum a_i,
 \qquad s_\lambda=m-l,
\]
let $G_\lambda=\prod_i\mathfrak S_{a_i}$ act on $S^l$, and set
\[
 e_\lambda=\frac1{|G_\lambda|}
       \sum_{g\in G_\lambda}[\Gamma_g]
 \in\operatorname{CH}^{2l}(S^l\times S^l)_{\mathbf Q}.
 \tag{HM20.2}\label{r20hm:averaging}
\]
This is an actual idempotent since multiplication in the
finite group ring gives $e_\lambda^2=e_\lambda$.

The primary theorem of de Cataldo--Migliorini gives
\[
 h(H)\simeq\bigoplus_{\lambda\vdash m}
       (h(S^{l(\lambda)}),e_\lambda)
                              \otimes\mathbb L^{s_\lambda}.
 \tag{HM20.3}\label{r20hm:dcm}
\]
The precise source is M. A. de Cataldo and L. Migliorini,
\emph{The Chow Groups and the Motive of the Hilbert Scheme
of Points on a Surface}, J. Algebra 251 (2002), 824--848,
Theorem 6.2.1, printed p.~846, and Remark 6.2.2; the latter
explicitly permits any ground field.  Proposition 6.1.5
provides the diagonal identity.  The primary author PDF is
\url{https://www.math.stonybrook.edu/~mde/MyPublishedPapers/MotiveHilbSchJournOfAlg.pdf}.
In particular this is an actual Chow-motive theorem over $K$,
not a theorem restricted to complex cohomology.

For clarity, its incidence cycle before quotienting is
\[
 \Gamma_\lambda=
 \bigl\{((x_1,\ldots,x_l),\xi):
       \pi(\xi)=\lambda_1x_1+\cdots+\lambda_lx_l\bigr\}_{\rm red}
 \subset S^l\times H,
 \tag{HM20.4}\label{r20hm:incidence}
\]
where $\pi:H\to S^{(m)}$ is Hilbert--Chow.  Its dimension
is $m+l$, so its codimension in $S^l\times H$ is $m+l$.
The theorem normalizes these cycles, their transposes and
the finite averages \eqref{r20hm:averaging} to give actual
maps $u_\lambda,v_\lambda$ with
\[
 \sum_\lambda v_\lambda\circ u_\lambda=\Delta_H,
 \quad
 u_{\lambda*}:\operatorname{CH}^r(H)\longrightarrow
                  \operatorname{CH}^{r-s_\lambda}(S^l),
 \quad
 v_{\lambda*}:\operatorname{CH}^{r-s_\lambda}(S^l)
                    \longrightarrow\operatorname{CH}^r(H).
 \tag{HM20.5}\label{r20hm:degrees}
\]
The degree calculation is $m+l-2m=-s_\lambda$ for the
transpose and $m+l-2l=s_\lambda$ for the forward map.
The maps in \eqref{r20hm:degrees} are the mutually inverse
maps of the cited Chow theorem, including its rational
normalization constants.  The singular quotient
$S^l/G_\lambda$ is represented by its projector on the
smooth power; it is never used as a smooth Green target.
Every construction is defined over $K$ and commutes with
field extension.

\subsection{The local and global full-domain consequence}

\begin{theorem}[Products of all point Hilbert schemes of arbitrary surfaces]
\label{r20hm:main}
Fix a finite valued $K/\mathbf Q_p$ with its original absolute value.
Let $S_1,\ldots,S_b/K$ be arbitrary smooth projective geometrically
integral surfaces, and let $m_i\geq1$.  Let $C_1,\ldots,C_c/K$
be arbitrary smooth projective geometrically integral curves.
Let $T/K$ be smooth projective geometrically integral with
potentially good projective reduction; a point is permitted.
Then
\[
 X=T\times\prod_{j=1}^{c}C_j\times
                    \prod_{i=1}^{b}S_i^{[m_i]}
 \tag{HM20.6}\label{r20hm:family}
\]
has strong G for every specified geometrically homologically
trivial rational cycle in every $1\leq r\leq\dim X$.
The conclusion holds over every finite extension of $K$ for
all its new geometrically homologically trivial cycles.
If $b=c=0$, the assertion is the frozen potentially-good case.
\end{theorem}

\begin{proof}
Tensor the finitely many actual identities
\eqref{r20hm:degrees} with each other and with the diagonals
of $T$ and the curves.  For every tuple of partitions
$\boldsymbol\lambda$, this yields maps through the smooth
projective variety
\[
 V_{\boldsymbol\lambda}
   =T\times\prod_j C_j\times\prod_i S_i^{l(\lambda_i)}
\]
with shift $s=\sum_i(m_i-l(\lambda_i))$, and their composites
sum to $\Delta_X$ in rational Chow.  These are actual
cycles and an actual rational-equivalence certificate.
DP20 proves full homological obstruction vanishing on
every $V_{\boldsymbol\lambda}$, using its allowed arbitrary
potentially-good coefficient.  The transfer argument above
then gives zero obstruction on $X$ in every degree, and
VT17 gives the desired exact-support actual Green solution.

To check the shift, write $n=\dim X$.  Then
$\dim V_{\boldsymbol\lambda}=n-2s$, and the image of a
codimension-$r$ cycle is in codimension $r-s$.  Homological
triviality is preserved.  If this degree is outside the
Chow range, or is zero, the image is zero as explained
above; otherwise it is in precisely DP20's asserted range.

For any finite $L/K$, base change the actual diagonal
identity and reapply DP20 to the varieties over $L$ and to
the new image cycles.  This does not assume they are old
pullbacks.  DP20's auxiliary alterations and field extensions
are internal to that theorem, with the frozen full-degree
actual descent.  The final VT17 construction remains on
the originally specified support over $L$, with all actual
compact tests and its prescribed absolute value.
\end{proof}

\begin{corollary}[Full complementary Beilinson--Bloch heights]
\label{r20hm:global}
Let $F$ be a number field.  Let the curves and surfaces in
\eqref{r20hm:family} be arbitrary smooth projective
geometrically integral varieties over $F$, and suppose that
$T/F$ has potentially good projective reduction at every
finite place.  For $n=\dim X$ and $1\leq r\leq n$, PH16
is defined on the full groups
\[
 \operatorname{CH}^{r}(X)_{\mathbf Q,\mathrm{geom\ hom}}
 \times
 \operatorname{CH}^{n+1-r}(X)_{\mathbf Q,\mathrm{geom\ hom}}
 \longrightarrow\mathbf R.
 \tag{HM20.7}\label{r20hm:height}
\]
The same holds after every finite number-field extension.
\end{corollary}

\begin{proof}
Apply the local theorem at every finite completion, and
after every finite extension of those completions.
All local obstructions on both full homological Chow groups
vanish.  PH16 and GA17 give the full intrinsic height
domain and the regular-model admissible Chow lifts.
No direct semistable model of the Hilbert scheme is required.
\end{proof}

\paragraph{A concrete middle range and the exact limit.}
For an arbitrary surface $S/F$, the sixfold $S^{[3]}$ has
the full pairing in codimensions $3$ and $4$.  Its three
partitions have lengths $3,2,1$ and shifts $0,1,2$, so the
correspondence images of the two arguments have respective
codimensions $(3,4)$ on $S^3$, $(2,3)$ on $S^2$ and $(1,2)$
on $S$.  In every row the sum is the target dimension plus
one.  This is a nonproduct geometry reached by the actual
diagonal identity.  Neither a cohomological motive splitting,
general finite-dimensionality of Chow motives, positivity
of heights, nor all-prime local-link comparison is inferred.
The general G and arbitrary-variety full-height problems
remain open.
